\section{Introduction}
Three-variable first-order logic ($\FOthree$) over binary relations and the
calculus of relations
(\CoR) define the same binary queries, an equivalence going back to
Tarski's 1940s work~\citep{tarski-givant1987}. \CoR\ is the algebra of
binary relations under union $P\cup Q$, intersection $P\cap Q$, complement
$\overline P$, converse $P^\smile=\set{(y,x):(x,y)\in P}$, and the
composition, or relative product,
$P;Q=\set{(x,y):\exists z.\,(x,z)\in P,\ (z,y)\in Q}$. \CoR\ has no
individual variables, so quantifier elimination has to go through the
composition:
\begin{equation}\label{eq:basic-identity}
 \exists z\,\bigl(P(x,z)\land Q(z,y)\bigr)\iff(P;Q)(x,y).
\end{equation}
Using \eqref{eq:basic-identity} together with bottom-up disjunctive normal
form (DNF)
conversion~\citep{hilbert-ackermann1928,kleene1952,schoening1989},
$\FOthree$ formulas can be translated into \CoR\ with a single-exponential
blow-up. Both were in place before Tarski took up the question
(see Section~\ref{sec:relwork}). For an example, fix $k\ge2$ and consider the
$\FOthree$ formula
\begin{equation}\label{eq:omega}
\Psi_k(x,y) := \exists z\,\bigwedge_{i=1}^k
       \bigl(R_i(x,z)\lor R_i(z,y)\bigr),
\end{equation}
which asks for a single witness $z$ satisfying $R_i(x,z)\lor R_i(z,y)$ for
every $i\in\{1,\dots,k\}$. To apply \eqref{eq:basic-identity}, the body
under the quantifier must be factored into the form
$\alpha(x,z)\land\beta(z,y)$. Thus, by DNF conversion, the body of
\eqref{eq:omega} expands into $2^k$ disjuncts, one for each subset
$I\subseteq[k]$, the disjunct for $I$ asserting that $R_i(x,z)$ holds for
$i\in I$ and $R_i(z,y)$ for $i\notin I$. Applying
\eqref{eq:basic-identity} to each disjunct gives an equivalent \CoR\ term of
size $2^{O(k)}$,
\begin{equation}\label{eq:omega-upper}
 \CoR_{\Psi_k} := \bigcup_{I\subseteq[k]}
 \biggl(\Bigl(\bigcap_{i\in I}R_i\Bigr);\Bigl(\bigcap_{i\notin I}R_i\Bigr)\biggr),
\end{equation}
where an empty intersection denotes the full relation $\one$, which holds at
every pair. For $k=3$ the eight choices of $I$ give eight disjuncts and
eight compositions:
\[
\renewcommand{\arraystretch}{1.15}
\begin{array}{c@{\qquad}l@{\qquad}l}
 I&\text{disjunct of the DNF}&\text{term}\\\hline
 \{1,2,3\}&R_1(x,z)\land R_2(x,z)\land R_3(x,z)&(R_1\cap R_2\cap R_3);\one\\
 \{1,2\}&R_1(x,z)\land R_2(x,z)\land R_3(z,y)&(R_1\cap R_2);R_3\\
 \{1,3\}&R_1(x,z)\land R_3(x,z)\land R_2(z,y)&(R_1\cap R_3);R_2\\
 \ \ \vdots&\qquad\qquad\vdots&\qquad\ \ \vdots\\
 \varnothing&R_1(z,y)\land R_2(z,y)\land R_3(z,y)&\one;(R_1\cap R_2\cap R_3)
\end{array}
\]
That the translation is exponential in general is classical.
\begin{proposition}[well known; see Section~\ref{sec:relwork}]\label{prop:upper-translation}
Every \FOthree\ formula $\varphi$ over binary relations with at most two
free variables is
equivalent, on all structures, to a \CoR\ term of size $2^{O(|\varphi|)}$,
where $|\cdot|$ counts the nodes of a syntax tree.
\end{proposition}

Whether this single-exponential upper bound is optimal has remained a
long-standing open question, and no matching lower bound for full \CoR\ has
been known since the correspondence itself was established.

In this paper we prove that the blow-up is unavoidable: the
$2^{\Omega(|\varphi|)}$ lower bound holds already over finite structures,
already for positive \FOthree\ formulas with a single quantifier, and even
when the \CoR\ side is a circuit rather than a term.

\subsection{Main results}\label{sec:main-results}
Write $|\varphi|$ and $|t|$ for the number of nodes in the syntax tree of a
formula $\varphi$ or a term $t$, and $\cost{t}$ for the number of
\emph{composition} nodes of $t$; in the second measure the Boolean
operations and converse are free. A \emph{relational circuit} $\Comp$ is a
term whose subterms may be shared, that is, a directed acyclic graph, and
$\gcost{\Comp}$ counts its distinct composition gates. Write
$t\equiv_{\rm fin}\varphi$ if $t$ and $\varphi$ define the same relation on
every finite structure, and likewise for circuits. The first hard query is
$\Psi_k$ itself, and the $2^k$ compositions of \eqref{eq:omega-upper} are
close to necessary.

\begin{theorem}[Disjunctive clauses]\label{main:omega}
Let $k\ge2$. Every \CoR\ term $t\equiv_{\rm fin}\Psi_k$ satisfies
\[
 \cost{t}\ \ge\ \Bigl\lceil\tfrac12\binom{k}{\lfloor k/2\rfloor}\Bigr\rceil
 \qquad\text{and}\qquad
 |t|\ \ge\ \binom{k}{\lfloor k/2\rfloor}+1,
\]
and the bound on $\cost{t}$ holds more generally for circuits: every
relational
circuit $\Comp\equiv_{\rm fin}\Psi_k$ has
$\gcost{\Comp}\ge\bigl\lceil\frac12\binom{k}{\lfloor k/2\rfloor}\bigr\rceil$.
\end{theorem}

Since $\binom{k}{\lfloor k/2\rfloor}=\Theta(2^k/\sqrt k)$, this falls short
of the $2^k$ compositions of \eqref{eq:omega-upper} by a factor of
$\Theta(\sqrt k)$. The second hard query closes the gap. It asks
the two pairs $(x,z)$ and $(z,y)$ to satisfy the
same one of two symbols in each clause, rather than to share out the $k$
clauses between them. For $k\ge1$
let
\begin{equation}\label{eq:phi}
\Phi_k(x,y) := \exists z\,\bigwedge_{i=1}^k
       \Bigl(\bigl(P_i(x,z)\land P_i(z,y)\bigr)\lor
             \bigl(Q_i(x,z)\land Q_i(z,y)\bigr)\Bigr),
\end{equation}
over the $2k$ independent symbols $P_1,Q_1,\ldots,P_k,Q_k$. Clause $i$ holds
at a witness $z$ when both edges $(x,z)$ and $(z,y)$ lie in $P_i$, or both
lie in $Q_i$.

\begin{theorem}[Profile agreement]\label{main:phi}
Let $k\ge1$. Every \CoR\ term $t\equiv_{\rm fin}\Phi_k$ satisfies
$\cost{t}\ge2^{k-1}$ and $|t|\ge2^k+1$, and the bound on $\cost{t}$ holds
more
generally for circuits: every relational circuit
$\Comp\equiv_{\rm fin}\Phi_k$ has $\gcost{\Comp}\ge2^{k-1}$. The explicit
term $\CoR_{\Phi_k}$ of \eqref{eq:phi-upper} uses $2^k$ compositions, so the
composition count
is determined within a factor of two.
\end{theorem}

Both formulas are positive, use a single quantifier, and have size
$\Theta(k)$ over a vocabulary that grows with $k$
(Remark~\ref{rem:vocabulary}), so the two theorems answer the question
above.

\begin{corollary}[Succinctness gap]\label{cor:gap}
The translation of Proposition~\ref{prop:upper-translation} is optimal up to
the constant in the exponent: there are positive \FOthree\ formulas
$\varphi$ with one quantifier, of arbitrarily large size, such that every
relational circuit equivalent to $\varphi$ over finite structures has size
$2^{\Omega(|\varphi|)}$.
\end{corollary}
\begin{proof}
Take $\varphi=\Psi_k$, of size $\Theta(k)$. By Theorem~\ref{main:omega} a
circuit for it has at least $\frac12\binom{k}{\lfloor k/2\rfloor}$
composition gates, and $\binom{k}{\lfloor k/2\rfloor}\ge2^k/(k+1)$, so its
size is at least $2^{k-O(\log k)}=2^{\Omega(|\varphi|)}$, while
\eqref{eq:omega-upper} is an equivalent term of size
$O(k2^k)=2^{O(|\varphi|)}$. Taking $\varphi=\Phi_k$
gives $2^k$ in place of $2^k/(k+1)$.
\end{proof}

Section~\ref{sec:overview} outlines the proof of
Theorem~\ref{main:omega}, which works also for Theorem~\ref{main:phi}.

What Corollary~\ref{cor:gap} leaves open is the same question over a
vocabulary \emph{fixed} in advance. The formulas above use $k$ and $2k$
relation symbols respectively,
and $|\varphi|=\Theta(k)$ because the size charges an occurrence of a symbol
one node. Over a single binary relation symbol we obtain, for positive formulas, a gap of
$2^{\Omega(n/\log n)}$ (Theorem~\ref{thm:af-fixed}); whether $2^{\Omega(n)}$
holds over a fixed vocabulary we do not know
(Remark~\ref{rem:vocabulary}).

\subsection{Additional results}\label{sec:additional-results}

\paragraph{Size-specific and randomized targets.}
A \CoR\ term is a single expression, evaluated on structures of every size,
so Theorems~\ref{main:omega} and~\ref{main:phi} are uniform lower bounds.
The same bounds hold non-uniformly, with the term or circuit chosen for each
domain size $n$, and a bounded probability of error costs only a constant
factor.

\begin{theorem}[informal; see Theorem~\ref{thm:fixed-size} and
Corollary~\ref{cor:randomized}]\label{main:robust}
Let $\varphi$ be $\Psi_k$ or $\Phi_k$. For every $n$ above a threshold
$n_0=2^{O(k)}$, the lower bounds of Theorems~\ref{main:omega}
and~\ref{main:phi} hold for every term or circuit that agrees with $\varphi$
on all structures with exactly $n$ elements, even when that term is chosen
depending on $n$. A randomized circuit that is correct with probability at
least $1-\varepsilon$ on every input, for some $\varepsilon<1/2$, still
needs a $(1-2\varepsilon)$ fraction of that same number of compositions, in
expectation.
\end{theorem}

\paragraph{Matrix query languages.}
The composition $P;Q$ can be replaced by $\bigvee_zP(x,z)\odot Q(z,y)$, with
$\odot$ an arbitrary binary operation on the scalars and $\bigvee$ an
arbitrary idempotent join (Section~\ref{sec:responses}), and the lower
bounds continue to hold; the min-plus matrix product
$\min_z\bigl(P(x,z)+Q(z,y)\bigr)$ is one instance. \citet{brijder2022}
showed that MATLANG defines the same binary queries as a positive
relational algebra on semiring-annotated relations whose intermediate
results have at most three attributes; their translation into MATLANG is
exponential, and they asked whether that blow-up is unavoidable.

\begin{theorem}[informal; see Theorem~\ref{main:matrix}, and
Corollary~\ref{cor:positive} for non-idempotent semirings]\label{main:matrix-informal}
Fix the min-plus semiring, or max-plus, or max-min. For every $k$ there is a
positive relational algebra expression $E_k$ of size $O(k)$, with binary
output, whose intermediate results have at most three attributes. Every
MATLANG computation equivalent to $E_k$ on $n\times n$ inputs uses at least
$2^k/3$ matrix products contracting the dimension $n$, even with arbitrary
uniform entrywise functions and shared intermediates; $2^k$ products
suffice.
\end{theorem}

\subsection{Related work}\label{sec:relwork}

\paragraph{The expressiveness question.}
For historical background, we refer to the surveys by \citet{pratt1992},
\citet{maddux1991}, and \citet{burris-legris}, the monograph of
\citet{brady2000}, and the historical notes of \citet{tarski-givant1987}.
\CoR\ is older than first-order logic. De Morgan introduced composition
and converse in his papers on the syllogism of 1850 and
1860~\citep{demorgan1856,demorgan1864}, Peirce developed the algebra
through the 1870s and 1880s~\citep{peirce1870,peirce1880}, and Schr\"oder
devoted an entire volume to it in 1895~\citep{schroeder1895}. In contrast,
first-order logic was formalized only in Hilbert's lectures of 1917--18 and
in Hilbert and Ackermann's 1928
book~\citep{hilbert-lectures,hilbert-ackermann1928,moore1988}. In the
terminology of the time, to \emph{condense} a quantified statement about
relations meant eliminating its quantifiers to express it within the
calculus~\citep[p.~233]{loewenheim1915}---the translation of
Proposition~\ref{prop:upper-translation} under its older name. Schr\"oder
claimed in 1895 that every such statement could be
condensed~\citep[p.~551]{schroeder1895}. Korselt found a counterexample,
published by L\"owenheim in 1915~\citep{loewenheim1915}, which lies
exactly on the variable boundary studied here: ``the domain has at most two
elements'' requires three variables (\FOthree) and can be condensed,
whereas ``the domain has at most three elements'' requires four (\FOfour) and
cannot. Statements in \FOfour\ were thus known to lie outside \CoR\
before first-order logic existed as a formal system, while the status of
\FOthree\ remained open.
Tarski's 1941 paper presented further
\FOfour\ statements outside \CoR,
observed that the calculus remained at essentially the same stage of development as forty-five years earlier,
and posed the problem of characterizing the
condensable sentences~\citep[pp.~74, 89]{tarski1941}.
According to Tarski and
Givant, Tarski obtained the \FOthree\ characterization in the early 1940s
and presented it at a 1945 Berkeley seminar. It was first published with full
proofs in their 1987 monograph~\citep[pp.~xvii, 88--89]{tarski-givant1987},
and Maddux provided a modern treatment in his monograph~\citep{maddux2006}.
This settled the expressiveness question and forms the basis for
Proposition~\ref{prop:upper-translation} and Corollary~\ref{cor:gap},
but the computational cost of this translation was not examined.

\paragraph{The succinctness question.}
Expressive equivalence, however, implies nothing about formula size. Tarski
and Givant defined their translation recursively, providing no bound on its
growth~\citep[\S3.9]{tarski-givant1987}. The folklore translation behind
Proposition~\ref{prop:upper-translation} applies bottom-up DNF conversion to
the quantifier-free Boolean shell of a formula---the propositional normal-form
theorem of \citet{hilbert-ackermann1928}, read as a conversion procedure by
\citet{kleene1952} and \citet{schoening1989}---and then applies
\eqref{eq:basic-identity} to each disjunct. This bottom-up conversion is
the sole source of the exponential blow-up. In contrast, the reverse
translation is linear: structural induction turns a \CoR\ term of size $m$
into an \FOthree\ formula of size $O(m)$, so the blow-up occurs in one
direction only. Whether this exponential blow-up is unavoidable has remained
open since the equivalence was established. The query $\Psi_k$ is due to
\citet{nakamura2022}, who showed that it has no subexponential translation
into \emph{positive} \CoR\ (the fragment without complement), leaving the
case of full \CoR\ open. Theorems~\ref{main:omega} and~\ref{main:phi} settle
this question.

\paragraph{Succinctness of logics and lower-bound methods.}
Succinctness is a standard measure for comparing logics of the same
expressive power~\citep{grohe-schweikardt2005}, and gaps of every size are
known between fragments of first-order, temporal, and modal
logics~\citep{etessami-vardi-wilke2002,adler-immerman2003,dawar2007,french2013}.
The closest result to ours concerns finite-variable fragments over a
single binary relation: on finite linear orders, \FOfour\ is
exponentially more succinct than \FOthree, while two- and three-variable
logics are polynomially equivalent~\citep{grohe-schweikardt2005}. Such lower
bounds are proved using the formula-size games of Adler and
Immerman~\citep{adler-immerman2003,hella-vilander2019} or, equivalently, by
tracking a potential function along syntax trees~\citep{grohe-schweikardt2005}.
This framework is indifferent to negation (which merely swaps structure pairs)
and charges a bounded potential increase for each quantifier (corresponding
to composition in \CoR). The core difficulty lies in finding suitable
structures and a potential, neither of which transfers from linear orders to
our setting, for two concrete reasons.

First, on linear orders there is no gap to find: by the lower-bound theorem of
Grohe and Schweikardt~\citep{grohe-schweikardt2005}, an \FOthree\ formula of
size $s$ can only tell in which order its points come and whether two of
them lie at distance exactly $d$ for some $d\le4s^2$, and each such test is
a \CoR\ term of size $O(d)$, namely a $d$-fold composition of
$\mathit{succ}=(<)\cap\overline{(<\,;<)}$. As a formula of size $s$ is a
Boolean combination of the $s^{O(1)}$ available tests, it has an equivalent
\CoR\ term of size $s^{O(1)}$; so \FOthree\ and \CoR\ terms are
polynomially equivalent there.

Second, the mechanism of their gap is fundamentally different. A fourth variable
permits the divide-and-conquer recursion familiar from Savitch's theorem~\citep{savitch1970},
\[
 \varphi_m(x,y) = \exists z\,\forall u\,\bigl((u=x\lor u=y)\to
 \varphi_{m-1}(z,u)\bigr),
\]
which expresses ``$x$ and $y$ are at distance $2^m$'' in size $O(m)$ by
using one copy of $\varphi_{m-1}$ for both halves of the
path~\citep{grohe-schweikardt2005}; with three variables the recursion must
write $\varphi_{m-1}$ twice at every level, and their lower bound shows
that no other \FOthree\ formula does better than $2^{\Omega(m)}$.
Reuse is exactly what a circuit provides: a relational circuit expresses
distance $2^m$ with $m$ compositions, by squaring $S_{i+1}=S_i;S_i$ from
$S_0=\mathit{succ}$, whereas a term needs size $2^{\Omega(m)}$ by their
bound and the linear translation of terms into \FOthree. Our bounds hold for
circuits---where on linear orders the gap actually runs in the opposite
direction---showing that the gap between \FOthree\ and \CoR\ is not simply
one of reuse.

This gap stems from how the two formalisms handle witnesses: in $\Psi_k$, the
quantifier binds $R_i(x,z)\lor R_i(z,y)$ for all $i$ simultaneously,
constraining both edges at once. In contrast, a composition $P;Q$ splits the
edges into $P$ and $Q$. The classical translation \eqref{eq:omega-upper} must
therefore spend one composition on each way of partitioning the $k$ clauses
between the two edges, an overhead that Theorem~\ref{main:omega} proves to be
unavoidable. \citet{nakamura2022} proved the exponential bound for positive
\CoR\ using a formula-size game; however, the argument relies fundamentally on
positivity and does not extend to complement, which exchanges the two sides
and thus requires a symmetric potential. Our preservation theorem supplies
one: for a single finite structure and the family of substructures obtained by
deleting one part each, the potential of a subterm is a set of parts outside
of which every deletion is invisible to it. Boolean operations and converse
contribute nothing to this set, while a composition contributes at most two
parts; because these parts combine by union rather than by addition, a shared
gate is charged only once, directly yielding the bound for circuits. Finally,
$\Psi_k$ and $\Phi_k$, suitably interpreted, detect every deletion.

\paragraph{Composition as matrix multiplication.}
Since composition corresponds to a Boolean matrix product, a \CoR\ term with $m$
compositions evaluates a query using $m$ fast matrix multiplications in time
$O(m\,n^{\omega})$, rather than the $O(n^3)$ time of naive enumeration,
where $\omega$ is the matrix multiplication exponent. Bounded-variable queries
can be evaluated in time polynomial in both the query and the
structure~\citep{vardi1995}. \citet{Williams2014} exploited this speedup for
three-quantifier graph properties, leading to a fine-grained complexity
theory of first-order queries parameterized by variable count and quantifier
structure~\citep{GaoImpagliazzoKolokolovaWilliams2017,BringmannFischerKunnemann2019,GaoImpagliazzo2019}.
The composition counts in Theorems~\ref{main:omega} and~\ref{main:phi}
therefore lower-bound the number of matrix products required by any
relational evaluation of the query, a bound stated directly in
Theorem~\ref{main:matrix}. These results are not lower bounds for general
Boolean circuits with access to individual input entries. In our cost measure,
all Boolean operations (including complement) are free. In contrast,
classical monotone lower bounds for a single Boolean matrix
product~\citep{kerr1970,pratt1975,paterson1975,mehlhorn-galil1976} and
exact bounds for straight-line programs over semirings~\citep{jerrum-snir1982}
count individual gates; see \citet{jukna2012} for a survey.

\paragraph{Relational and matrix query languages.}
Algebraic representations of logical queries have been central to query
processing since Codd~\citep{Codd1970,Codd1972}. Binary relation algebra and its
fragments form the algebraic core of graph and navigational query
languages~\citep{FletcherEtAl2015,HellingsEtAl2021,HellingsEtAl2022}. In these
settings, logic-to-algebra translations can incur exponential
blowups~\citep{BaranyTenCateOtto2012}. The succinctness of relational
representations is therefore a central concern~\citep{FletcherEtAl2015},
because the size of the translated expression directly dictates the cost of
query optimization and evaluation. MATLANG is a matrix query language built
from matrix multiplication, transpose, ones-vectors, diagonalization, and
pointwise functions~\citep{brijder2019tods}; see \citet{geerts2021} for a
survey. For relations annotated with values in a semiring, in the sense of
provenance semirings~\citep{green2007},
\citet{brijder2022} proved that binary-output ARA(3)---the three-attribute
annotated relational algebra---is expressively equivalent to MATLANG,
mirroring the equivalence between \FOthree\ and \CoR. An earlier manuscript
asked whether this exponential arity elimination (and by analogy, the
classical \FOthree\ translation) is unavoidable~\citep[p.~9]{brijder2019}.
Theorem~\ref{main:matrix} resolves this question affirmatively for
commutative idempotent semirings of finite join breadth, while
Section~\ref{semi:positive-section} extends the bound to additional semirings
under a restricted operation set.


\subsection{Technical overview}\label{sec:overview}
We outline the proof of Theorem~\ref{main:omega}. Theorem~\ref{main:phi}
uses the same construction and the same preservation theorem, and differs
only in the interpretation of the relation symbols.

\subsubsection{Outline}\label{ov:outline}
The term \eqref{eq:omega-upper} contains one composition for each way of
distributing the $k$ clauses between the two edges of the witness. A more
economical term might use complementation, factor out common subterms, or
share them, and a normal-form argument would have to account for all of
these possibilities. We therefore do not analyze the syntax of a candidate
term at all. Instead, we fix a single finite structure, delete parts of it,
and ask which of the resulting substructures a given term can still
distinguish from the original.

We say that a term or a query \emph{distinguishes} two structures over the
same vocabulary, one of them a substructure of the other, if the relations
that it defines on them differ at some pair of vertices of the smaller
structure; equivalently, the term or query \emph{detects} the deletion that
produced the smaller structure. The proof has three steps.
\begin{enumerate}[nosep,label=(\roman*)]
\item We build one finite structure $\M$ (Figure~\ref{fig:structure}).
 Apart from a hub and some anchors, its vertices are partitioned into
 \emph{cells} $\cell\lambda$, one for each color $\lambda$ in a finite set
 $\Col$ of at least two colors. The only substructures we use are those
 obtained by deleting entire cells: for nonempty $J\subseteq\Col$, let
 $\M\res J$ be $\M$ with every cell $\cell\lambda$ such that
 $\lambda\notin J$ removed. Thus no vertex is ever removed individually,
 and the hub and the anchors survive every deletion.
\item To a \CoR\ term $t$ with $m$ compositions we attach a set
 $K\subseteq\Col$ of at most $2m$ colors, and show that $t$ does not
 distinguish $\M$ from $\M\res J$ for any nonempty $J\supseteq K$
 (Theorem~\ref{thm:preservation}): deleting cells is invisible to $t$ as
 long as every cell whose color lies in $K$ is retained.
\item We interpret the relation symbols of $\Psi_k$ on $\M$ so that
 $\Psi_k$ does distinguish $\M$ from $\M\res(\Col\setminus\{\lambda\})$, for
 every single color $\lambda$ (Section~\ref{sec:logical}).
\end{enumerate}
Suppose $t\equiv_{\rm fin}\Psi_k$, and let $\lambda$ be any color. By~(iii)
the formula distinguishes $\M$ from $\M\res(\Col\setminus\{\lambda\})$, and
hence so does $t$. Since $|\Col|\ge2$, the set $\Col\setminus\{\lambda\}$ is
nonempty, so by~(ii) this is possible only if
$\Col\setminus\{\lambda\}\not\supseteq K$, that is, only if $\lambda\in K$.
Hence $K=\Col$ and $2m\ge|\Col|$. Taking for $\Col$ the $\lfloor
k/2\rfloor$-element subsets of $[k]$ gives the composition bound of
Theorem~\ref{main:omega}; the bound on $|t|$ follows from
\eqref{eq:tree-size}.

The remainder of this section treats the three steps in order.
Section~\ref{ov:structure} constructs $\M$ and specifies the pairs of
vertices at which a query is read off; Sections~\ref{ov:types}
to~\ref{ov:wholeterm} prove~(ii) by showing that what a term computes on
$\M$ is too coarse to detect an individual cell; and
Section~\ref{ov:formula} establishes~(iii).

\begin{figure}[t]
\centering
\begin{tikzpicture}[
  font=\small,
  cellnode/.style={draw,ellipse,minimum width=2.1cm,minimum height=1.0cm,thick},
  vtx/.style={circle,fill=black,inner sep=1.1pt},
  anch/.style={draw,rectangle,thick,minimum size=6mm,fill=white,inner sep=1pt}
]
 % ---------- hub ----------
 \node[vtx,label={[label distance=1pt]above:{hub $h$}}] (hub) at (0.4,3.6) {};
 % ---------- cells ----------
 \node[cellnode] (c1) at (-3.0,0) {};
 \node[cellnode] (c2) at ( 0.4,0) {};
 \node[cellnode] (c3) at ( 3.8,0) {};
 \node[below=2pt of c1,font=\footnotesize] {$\cell{\lambda_1}$};
 \node[below=2pt of c2,font=\footnotesize] {$\cell{\lambda_2}$};
 \node[below=2pt of c3,font=\footnotesize] {$\cell{\lambda_3}$};
 \foreach \c in {1,2}{
   \node[vtx] (c\c a) at ($(c\c)+(-0.58,0.12)$) {};
   \node[vtx] (c\c b) at ($(c\c)+(0.0,-0.22)$) {};
   \node[vtx] (c\c c) at ($(c\c)+(0.58,0.12)$) {};
 }
 % four dots in the third cell: the number drawn carries no meaning
 \node[vtx] (c3a) at ($(c3)+(-0.68,0.08)$) {};
 \node[vtx] (c3b) at ($(c3)+(-0.23,-0.24)$) {};
 \node[vtx] (c3c) at ($(c3)+(0.23,0.18)$) {};
 \node[vtx] (c3d) at ($(c3)+(0.68,-0.12)$) {};
 % ---------- anchors: one per color ----------
 \node[anch] (x1) at (-6.7,3.6) {$x_{\lambda_1}$};
 \node[anch] (x2) at (-6.7,2.5) {$x_{\lambda_2}$};
 \node[anch] (x3) at (-6.7,1.4) {$x_{\lambda_3}$};
 \node[font=\footnotesize,align=center] at (-6.7,4.45) {anchors:\\one per color};
 % anchor edges carry no color; only a few are drawn
 \draw[lightgray,thin] (x2) -- (hub);
 \draw[lightgray,thin] (x3) -- (hub);
 \draw[lightgray,thin] (x1) -- (x2);
 \draw[lightgray,thin] (x2) -- (x3);
 \foreach \t in {c1,c2,c3}{ \draw[lightgray,thin] (x2) -- (\t.north west); }
 % ---------- the read-off pair ----------
 \draw[very thick] (x1) -- (hub);
 \node[font=\footnotesize,fill=white,inner sep=1.5pt] at (-3.15,3.83)
      {read-off pair $(x_{\lambda_1},h)$};
 % ---------- hub edges: one line style per color ----------
 \foreach \d in {a,b,c}{ \draw[thick]         (hub) -- (c1\d); }
 \foreach \d in {a,b,c}{ \draw[thick,dashed]  (hub) -- (c2\d); }
 \foreach \d in {a,b,c,d}{ \draw[thick,dotted]  (hub) -- (c3\d); }
 \node[font=\footnotesize,fill=white,inner sep=1pt] at (-1.54,1.60) {$\lambda_1$};
 \node[font=\footnotesize,fill=white,inner sep=1pt] at ( 0.40,1.60) {$\lambda_2$};
 \node[font=\footnotesize,fill=white,inner sep=1pt] at ( 2.34,1.60) {$\lambda_3$};
 % ---------- random colors among cell vertices ----------
 \draw[gray] (c1c) -- (c2a);
 \draw[gray] (c2c) -- (c3a);
 \draw[gray] (c1b) to[bend right=30] (c3b);
 \draw[gray] (c2a) to[bend right=45] (c2c);
 \node[gray,font=\footnotesize] at (2.15,-1.92)
      {random colors, inside and between cells};
 % ---------- the deletion ----------
 \draw[thick,densely dashed,rounded corners=3pt]
      ($(c1)+(-1.40,0.72)$) rectangle ($(c1)+(1.40,-1.08)$);
 \node[font=\footnotesize,align=center] at (-3.0,-1.92)
      {deleted in\\$\M\res(\Col\setminus\{\lambda_1\})$};
\end{tikzpicture}
\caption{The structure $\M$ with three colors. Each ellipse is a cell and
the dots inside it are cell vertices; all cells have the same (large) number
of vertices, and the number drawn here is not significant. Every pair of
\emph{distinct} vertices is an edge; only a few are drawn. The hub is joined
to each vertex of $\cell\lambda$ by an edge of color $\lambda$ (solid,
dashed and dotted for $\lambda_1,\lambda_2,\lambda_3$), so the hub alone can
distinguish the cells. Edges between cell vertices, inside a cell or across
cells, carry random colors (gray) and realize every color between every
pair of cells. Anchors, one per color, are joined to all vertices and their
edges carry no color (light gray); they are distinguished by the relation
symbols and are never deleted. Deleting the cell $\cell{\lambda_1}$ removes
all of its vertices but leaves the read-off pair $(x_{\lambda_1},h)$ in
place.}
\label{fig:structure}
\end{figure}

\subsubsection{The structure (Section~\ref{sec:structure})}\label{ov:structure}
Recall that $\Col$ is a finite set of \emph{colors}. The vertices of $\M$
are of three kinds, each with a distinct role (Figure~\ref{fig:structure}):
\begin{center}
\renewcommand{\arraystretch}{1.15}
\begin{tabular}{@{}lll@{}}
\toprule
 &Number&Role\\
\midrule
hub $h$&one&the only vertex that can distinguish the cells\\
cells $\cell\lambda$&one per color $\lambda\in\Col$&contain the witnesses; a deletion removes a whole cell\\
anchors $x_\lambda$&one per color, never deleted&the query is read off at the pair $(x_\lambda,h)$\\
\bottomrule
\end{tabular}
\end{center}
Every pair of distinct vertices is an \emph{edge}, and every edge not
incident to an anchor carries a color $\colr(u,v)=\colr(v,u)\in\Col$. The
hub's edges announce the cell,
\[
 \colr(h,u)=\lambda\qquad\text{for }u\in\cell\lambda,
\]
while the edges between two \emph{cell vertices}, inside one cell or across
two, are colored at random. Lemma~\ref{lem:palette} then yields the
following \emph{extension property}:
\begin{quote}
for all \emph{distinct} cell vertices $u,v$, all colors $\mu,\nu$, and every
cell $\cell\lambda$, some $z\in\cell\lambda$ other than $u$ and $v$
satisfies $\colr(u,z)=\mu$ and $\colr(z,v)=\nu$.
\end{quote}
Thus a cell vertex sees every color in every cell and cannot distinguish one
cell from another, whereas the hub can: it sees each cell in a single color.
Edges incident to an anchor are left uncolored; anchors are instead
distinguished by the interpretation of the relation symbols
(Section~\ref{sec:logical}).


\paragraph{Why anchors are needed.}
We call a pair $(x_\lambda,h)$ a \emph{read-off pair}. Such a pair must
serve two purposes, one supplied by each of its two vertices.

First, it must survive every deletion. The test for a color $\lambda$
compares $\M$ with $\M\res(\Col\setminus\{\lambda\})$, so the query must be
read off at a pair that is present in both structures, and
Section~\ref{sec:logical} assigns a different pair to each color. A cell
vertex does not survive the deletion of its own cell, and the hub is a
single vertex, so the construction needs $|\Col|$ further vertices that are
never deleted and that the relation symbols can distinguish. These are the
anchors. They are indispensable rather than merely convenient: both hard
formulas have the form $\exists z\,\psi$ in which every atom of $\psi$
relates $z$ to $x$ or to $y$, and for formulas of this form
Lemma~\ref{lem:where} in Appendix~\ref{app:optimality} bounds the number
of detectable colors by $2|X|+2$, where $X$ is the set of anchors, for every
interpretation of the kind used in this paper (Section~\ref{ov:types}). The number of anchors is thus
precisely the resource that the lower bound consumes.

Second, it must be able to see the cell of a witness. The edge from a
witness $z\in\cell\lambda$ to the hub has color $\lambda$, whereas the edges
from a cell vertex into $\cell\lambda$ realize every color, by the extension
property. Hence the hub is the only vertex from which the cell of a witness
can be read off.

An anchor, by contrast, has no colored edges at all. The only color
associated with $x_\lambda$ is the $\lambda$ in its name; we call it the
\emph{label} of the anchor, and the interpretation of the relation symbols
in Section~\ref{sec:logical} will make it readable. Labels and colors are
the same objects; the two words merely indicate whether the object is
carried by an anchor or by an edge. Consequently, a query can compare a
witness against a label only at a read-off pair $(x_\lambda,h)$: the anchor
supplies the label, the hub supplies the color of the witness, and the query
tests whether the two agree.

This completes step~(i) of the outline. The read-off pairs play no further
role until Section~\ref{ov:formula}, where $\Psi_k$ performs exactly the
comparison just described and \eqref{eq:omega-decoder} records the result;
step~(ii), proved in Sections~\ref{ov:types} to~\ref{ov:wholeterm},
concerns all pairs of $\M$ simultaneously, and only step~(iii) uses one
designated pair per color.

\subsubsection{Types and type-constant relations (Section~\ref{sec:types})}\label{ov:types}
Pairs of vertices are sorted into finitely many \emph{types}, each of which
records only part of the information that the construction attaches to a
pair.
\needspace{10\baselineskip}
Concretely, the type of an ordered pair $(u,v)$ is
\begin{center}
\renewcommand{\arraystretch}{1.1}
\begin{tabular}{@{}ll@{}}
\toprule
the pair $(u,v)$ itself&if $u$ and $v$ are both anchors,\\
\emph{out of $x$}&if $u$ is the anchor $x$ and $v$ is not an anchor,\\
\emph{into $x$}&if $v$ is the anchor $x$ and $u$ is not an anchor,\\
\emph{loop}&if $u=v$ and $u$ is not an anchor,\\
\emph{color $\lambda$}&if $u\ne v$, neither is an anchor and $\colr(u,v)=\lambda$.\\
\bottomrule
\end{tabular}
\end{center}
We write $\tp(u,v)$ for this type, and $\ctype\lambda$ for the type
\emph{color $\lambda$}; the formal definition, with names for the remaining
types, is \eqref{eq:types}. The number of types is finite and independent
of the size of the cells: there is one type per ordered pair of anchors, one
\emph{out of} and one \emph{into} type per anchor, one type for the loops,
and one type per color.

A type records which anchors are involved, whether the two endpoints
coincide, and the color of the edge between them. It does not record which
cell an endpoint lies in, which vertex of that cell it is, or whether the
hub is an endpoint at all: in Figure~\ref{fig:structure}, an edge inside
$\cell{\lambda_2}$, an edge from $\cell{\lambda_1}$ to $\cell{\lambda_3}$,
and an edge from the hub into $\cell{\lambda_1}$ all have the same type
whenever they carry the same color. Types are thus a deliberate coarsening
rather than the finest available classification of pairs. In particular,
although the hub is clearly visible in $\M$, what makes it special is not
that types can see it but that the colors of its edges are deterministic
(Section~\ref{ov:twocolors}).

Call a binary relation $P$ on the vertices \emph{type-constant} if $P(u,v)$
depends only on $\tp(u,v)$, and write $P(\tau)$ for the common value of $P$
on the pairs of a type $\tau$. Not every relation is type-constant.
Consider, for example, the relation
$P=\cell{\lambda_1}\times\cell{\lambda_1}$, which holds at $(u,v)$ exactly
when both endpoints lie in the first cell of Figure~\ref{fig:structure}. It
holds at the loop $(u,u)$ of a vertex $u\in\cell{\lambda_1}$ but fails at
the loop $(h,h)$ of the hub, although these two pairs have the same type,
namely \emph{loop}; hence $P$ is not constant on that type. The relation
symbols of $\Psi_k$ and $\Phi_k$ will be interpreted by type-constant
relations (Section~\ref{sec:logical}), and type-constancy then propagates
to everything computed from them.

\paragraph{Everything computed is type-constant.}
Let $P,Q$ be type-constant. In
\[
 (P;Q)(u,v)=\bigvee_{z}P(u,z)\land Q(z,v),
 \qquad\text{the join taken over all vertices }z,
\]
the term that a witness $z$ contributes depends only on the pair of types
$\bigl(\tp(u,z),\tp(z,v)\bigr)$, and on nothing else about $z$. We call
this pair the \emph{profile} of $z$ at $(u,v)$. The value of the
composition is therefore determined by the \emph{set} of profiles that some
witness realizes: not by how many witnesses realize a profile, and not by
which vertices they are.

Lemma~\ref{lem:profile} states that this set of realized profiles depends
only on the type of $(u,v)$. If $(u,v)$ is an edge of color $\lambda$, for
instance, the realized profiles are (color $\mu$, color $\nu$) for every
pair of colors, by the extension property; (loop, color $\lambda$) and
(color $\lambda$, loop), realized by $z=u$ and $z=v$; and (into $x$, out of
$x$) for each anchor $x$. Hence $P;Q$ is again type-constant. Boolean
operations and converse trivially preserve type-constancy
(Corollary~\ref{cor:closure}), so every relation that a \CoR\ term computes
on $\M$ from type-constant inputs is a table with one entry per type. In
particular, no term can single out a cell, let alone an individual vertex.

\subsubsection{A composition sees at most two colors (Section~\ref{sec:responses})}\label{ov:twocolors}
Fix two type-constant relations $P,Q$ and abbreviate their values on the edges of
color $\lambda$ as
\[
 p_\lambda=P(\ctype\lambda),\qquad q_\lambda=Q(\ctype\lambda).
\]

As in step~(i) of the outline, delete the cells of all colors outside a
nonempty set $J\subseteq\Col$. A deletion removes witnesses, and a witness
matters only through its profile, so the deletion can change the value of
$P;Q$ at a pair only if it removes the \emph{last} witness of some profile
at that pair. At a pair with no hub endpoint this never happens: by the
extension property, a single surviving cell already realizes every profile
that the deleted cells realized. At a pair with a hub endpoint it can
happen, but only in a limited way. Every edge from the hub into
$\cell\lambda$ has color $\lambda$, so the hub sees a witness
$z\in\cell\lambda$ only through $p_\lambda$ and $q_\lambda$:
$P(h,z)=p_\lambda$ and $Q(z,h)=q_\lambda$. Consider the contribution of such
a $z$ to the join $(P;Q)(u,v)=\bigvee_zP(u,z)\land Q(z,v)$ at a pair $(u,v)$
that has the hub as an endpoint. The hub may be the left endpoint, the right
endpoint, or both, and in these three cases the contribution of $z$ is
\[
 \begin{aligned}
  p_\lambda\land Q(\tp(z,v))&\qquad\text{when }u=h\text{ and }v\ne h,\\
  P(\tp(u,z))\land q_\lambda&\qquad\text{when }v=h\text{ and }u\ne h,\\
  p_\lambda\land q_\lambda&\qquad\text{when }u=v=h.
 \end{aligned}
\]
At a pair of $\M\res J$, the anchors and the two endpoints themselves
survive the deletion and contribute the same terms regardless of $J$. The
only part of the join that $J$ can affect is therefore the part ranging over
the cell witnesses, that is, over the colors $\lambda\in J$ and, for each
such color, over the vertices of $\cell\lambda$. In the first case, as $z$
ranges over a cell, the factor $Q(\tp(z,v))$ takes the same set of values
whichever cell is chosen, by the extension property; the join therefore
factors as $\bigl(\bigvee_{\lambda\in J}p_\lambda\bigr)\land c$ with $c$
independent of $J$, and the only information about $J$ that survives is
whether some $\lambda\in J$ has $p_\lambda=1$. The second case is symmetric.
In the third case both edges of $z$ carry the color $\lambda$, and what
survives is whether some $\lambda\in J$ has $p_\lambda=q_\lambda=1$. Thus
$J$ can influence the value only through three bits (Lemma~\ref{lem:local}):
\begin{center}
whether some $\lambda\in J$ has $p_\lambda=1$;\quad whether some has
$q_\lambda=1$;\quad whether some has both.
\end{center}
Two colors always suffice to fix all three bits (Lemma~\ref{lem:and-width}).
If some color has $p_\lambda=q_\lambda=1$, that color alone fixes all three;
otherwise the third bit is $0$, and one color for each of the first two bits
suffices. We call such a set of at most two colors a \emph{support} of the
composition. As long as no cell whose color lies in the support is deleted,
the value of the composition on the surviving vertices is unchanged.

For example, Figure~\ref{fig:structure} has three colors, so a family
$(p_\lambda,q_\lambda)_{\lambda\in\Col}$ consists of two triples. Suppose
first that
\[
 (p_{\lambda_1},p_{\lambda_2},p_{\lambda_3})=(1,1,0),
 \qquad
 (q_{\lambda_1},q_{\lambda_2},q_{\lambda_3})=(0,1,1),
\]
that is, $P$ holds on the edges of colors $\lambda_1$ and $\lambda_2$, and
$Q$ holds on those of colors $\lambda_2$ and $\lambda_3$. Since
$p_{\lambda_2}=q_{\lambda_2}=1$, the color $\lambda_2$ alone fixes all three
bits, so $\{\lambda_2\}$ is a support, and both $\cell{\lambda_1}$ and
$\cell{\lambda_3}$ can be deleted without affecting the composition. Suppose
instead that
\[
 (p_{\lambda_1},p_{\lambda_2},p_{\lambda_3})=(1,0,0),
 \qquad
 (q_{\lambda_1},q_{\lambda_2},q_{\lambda_3})=(0,1,0).
\]
Now no color has both of its values equal to $1$: keeping $\lambda_1$ alone
would change the second bit and keeping $\lambda_2$ alone would change the
first, so every support must contain both colors, and
$\{\lambda_1,\lambda_2\}$ is one. Lemma~\ref{lem:and-width} states that this
is the worst case.

\subsubsection{From one composition to a whole term (Section~\ref{sec:simultaneous})}\label{ov:wholeterm}
A support controls a single composition. For a term $t$ with $m$
compositions, let $K$ be the union of the supports of all of its
compositions, so that $|K|\le2m$, and fix any nonempty $J\supseteq K$.
Induction along the evaluation order then shows that the value of every
subterm of $t$ on $\M\res J$ is the restriction of its value on $\M$
(Theorem~\ref{thm:preservation}): the Boolean operations and converse
commute with restriction, and at a composition the two operands already
agree by induction, while $J$ contains the support of that composition, so
Section~\ref{ov:twocolors} applies. Hence $t$ does not distinguish $\M$ from
$\M\res J$ for any nonempty $J\supseteq K$, which is step~(ii) of the
outline. A shared subterm is charged only once, so the same argument bounds
the number of composition gates of a circuit.

At a composition gate, the two relations being compared are always related
by an inclusion in one direction, for a trivial reason: deleting cells
removes witnesses and never adds any, so a pair related by the composition
on $\M\res J$ is related by it on $\M$ as well. Writing $D$ for the vertex
set of $\M$, $D_J$ for that of $\M\res J$, and $P\res D_J$ for the set of
pairs of a relation $P$ that have both endpoints in $D_J$,
\[
 (P\res D_J)\,;(Q\res D_J)\ \subseteq\ (P;Q)\res D_J
\]
holds for every $J$. What can fail is the reverse inclusion: a pair may be
related on $\M$ only through a witness that the deletion removed. Excluding
this is precisely the role of a support, which is why
Section~\ref{ov:twocolors} yields an \emph{equality} and not merely this
inclusion.

Complement is the operation for which the distinction between inclusion and
equality matters. The issue is not that complement fails to commute with
restriction: a subterm is complemented within the structure on which it is
evaluated, relative to $D_J^2$ on $\M\res J$ and to $D^2$ on $\M$, and
$D_J^2\setminus(P\res D_J)=(D^2\setminus P)\res D_J$ holds exactly, for
every relation $P$. The issue is that complement reverses inclusions rather
than preserving them. If only the inclusion displayed above were available,
the two sides for the complemented subterm would satisfy
\[
 D_J^2\setminus\bigl[(P\res D_J)\,;(Q\res D_J)\bigr]
 \ \supseteq\
 D_J^2\setminus\bigl[(P;Q)\res D_J\bigr],
\]
with the inclusion now pointing in the opposite direction, and a further
complement would reverse it again. No single inclusion can therefore be
maintained through a term whose complements occur at arbitrary depths. An
equality, by contrast, has no direction to reverse: $A=B$ yields
$D_J^2\setminus A=D_J^2\setminus B$, and union, intersection and converse
preserve equalities equally well. This is why complement gates require no
separate treatment, why no monotonicity is used anywhere in the argument,
and why the target may be an arbitrary \CoR\ term rather than a positive
one.

\subsubsection{The formula detects every deletion (Section~\ref{sec:logical})}\label{ov:formula}
It remains to interpret the $R_i$. Recall that $\Psi_k$ has one quantifier,
$\Psi_k(x,y)=\exists z\bigwedge_{i=1}^k\bigl(R_i(x,z)\lor R_i(z,y)\bigr)$.
Evaluated at a pair $(u,v)$ of vertices, it asks for a \emph{witness}: a
vertex $z$ such that every one of the $k$ clauses holds, either through the
edge $(u,z)$ or through the edge $(z,v)$.

So far $\Col$ has been an arbitrary finite set of colors. We now fix it: let
$\Col$ be the set of $\lfloor k/2\rfloor$-element subsets of $[k]$, and
write $A,B$ for colors in place of $\lambda,\mu$. There is one anchor $x_A$
for each color $A$.

We now define the relations $R_i$. For all colors $A$ and $B$,
\begin{center}
\begin{tabular}{@{}l@{\quad}l@{}}
$R_i(x_A,u)$ holds exactly when $i\in A$,&for every vertex $u$ that is not
 an anchor;\\
$R_i(u,v)$ holds exactly when $i\notin B$,&for every two distinct non-anchor
 vertices\\
&$u,v$ whose edge has the color $B$;\\
\end{tabular}
\end{center}
and $R_i$ is false on every remaining pair. This interpretation is
type-constant, as required by Section~\ref{ov:types}: the pairs from $x_A$
to a non-anchor vertex form a single type (edges at an anchor carry no
color), and the pairs of color $B$ form a single type, whether they run
from the hub into $\cell B$ or between two cell vertices.

At the read-off pair $(x_A,h)$, a
vertex $z$ of $\cell B$ satisfies clause $i$ if $i\in A$, through the edge
$(x_A,z)$, or if $i\notin B$, through the edge $(z,h)$; so it satisfies
all $k$ clauses exactly when $B\subseteq A$, and since $A$ and $B$ have the
same size, exactly when $B=A$. The hub is not a witness: $R_i$ is false on
the loop $(h,h)$, so at $z=h$ clause $i$ would need $i\in A$, and $A$
contains only $\lfloor k/2\rfloor$ of the $k$ indices. Nor is an anchor
$x_B$: $R_i$ is false on the pair $(x_A,x_B)$, so clause $i$ would need
$i\in B$, and $|B|=\lfloor k/2\rfloor$ as well. The witnesses at $(x_A,h)$
are therefore
exactly the vertices of the single cell $\cell A$, and
\[
 \Psi_k\text{ holds at }(x_A,h)\text{ in }\M\res J
 \iff A\in J,
\]
which is \eqref{eq:omega-decoder}. Deleting $\cell A$ therefore changes the
value of $\Psi_k$ at a pair that survives the deletion, so every color is
detectable, which completes step~(iii) of the outline.

The choice of labels is forced by the method, in a sense that can be made
precise. At a read-off pair, the formula compares the set announced by the
anchor with the set carried by the witness's edge to the hub, and a color is
detected precisely when exactly one color passes this comparison.
Recording, for each of the two sides, the set of clauses that it leaves
unsatisfied turns the comparison into a disjointness test between two
subsets of $[k]$, and detection into a cross-intersecting condition on
pairs of such subsets. Bollob\'as's set-pair inequality
(Lemma~\ref{lem:bollobas}), which bounds the size of such a family of pairs
of subsets of $[k]$ by $\binom{k}{\lfloor k/2\rfloor}$, then bounds the
number of detectable colors by the same quantity
(Proposition~\ref{prop:middle-layer} in Appendix~\ref{app:optimality}), and
the interpretation above attains this bound. This explains why
Theorem~\ref{main:omega} involves a central binomial coefficient rather
than $2^k$. The resulting gap of a factor
$\Theta(\sqrt k)$ against the $2^k$ compositions of \eqref{eq:omega-upper}
is not closed in this paper; by the same proposition, it is inherent to
$\Psi_k$ and to this construction, and not an artifact of the particular
relations we chose. The formula $\Phi_k$ avoids this loss by requiring
agreement rather than inclusion, which allows all $2^k$ subsets of $[k]$ to
serve as labels (Section~\ref{sec:phi}). Finally, the only property of
conjunction used in the argument is the one isolated in
Section~\ref{ov:twocolors}, namely that two colors fix its three bits.
Section~\ref{sec:matrix} replaces conjunction by a semiring product, and the
same argument goes through with three colors in place of two.

\subsection{Organization}
Section~\ref{sec:aggregation} fixes the conventions. The body then falls
into two parts, matching the two groups of results above. Part~I proves
Theorems~\ref{main:omega} and~\ref{main:phi}:
Section~\ref{sec:preservation} builds the structure $\M$ and proves the
preservation theorem, and Section~\ref{sec:logical} interprets the two
formulas on $\M$. Part~II collects the extensions:
Section~\ref{sec:robustness} passes to size-specific and randomized targets,
and Section~\ref{sec:matrix} treats semiring products and MATLANG. The two
can be read in either order, except that the fixed-size statements of
Section~\ref{sec:matrix} rest on Proposition~\ref{prop:fixed-test}.
Appendix~\ref{app:explicit} makes the coloring of $\M$ explicit,
Appendix~\ref{app:optimality} shows that the interpretations of
Section~\ref{sec:logical} are optimal, Appendix~\ref{app:kernels}
collects remarks on other kernels and on wider scalars, and
Appendix~\ref{app:anchor-free} proves the bounds of
Theorems~\ref{main:omega} and~\ref{main:phi}, and of
Section~\ref{sec:matrix}, without anchors for variants of the hard
queries.
